% 天体物理学（综述）
% license CCBYSA3
% type Wiki

本文根据 CC-BY-SA 协议转载翻译自维基百科\href{https://en.wikipedia.org/wiki/Astrophysics?utm_source=chatgpt.com}{相关文章}。


天体物理学是一门利用物理学与化学的方法和原理来研究天体及其现象的科学。\(^\text{[2]}\)作为该学科的奠基人之一，詹姆斯·基勒（James Keeler）曾指出，天体物理学“旨在探明天体的本质，而非它们在空间中的位置或运动——研究它们是什么，而非它们在哪里”，\(^\text{[3]}\)这一区别正体现于天体力学的研究范畴中。

天体物理学的研究对象包括太阳（太阳物理学）、其他恒星、星系、系外行星、星际介质以及宇宙微波背景辐射。\(^\text{[4][5]}\)科研人员在整个电磁波谱范围内观测这些天体的辐射，分析的物理特性包括光度、密度、温度及化学成分等。由于天体物理学的研究范围极其广泛，天体物理学家在研究中会综合运用多个物理学分支的概念与方法，包括经典力学、电磁学、统计力学、热力学、量子力学、相对论、核物理与粒子物理，以及原子与分子物理等。

在实践中，现代天文学研究通常涉及大量理论与观测物理学工作。天体物理学的主要研究领域包括：暗物质、暗能量、黑洞及其他天体的性质，以及宇宙的起源与最终命运。\(^\text{[4]}\)
理论天体物理学还研究多个主题，例如：太阳系的形成与演化、恒星动力学与演化、星系的形成与演化、磁流体力学、宇宙大尺度物质结构的形成、宇宙射线的起源、广义相对论与狭义相对论、量子宇宙学与物理宇宙学（即对宇宙最大尺度结构的物理研究），其中也包括弦论宇宙学与天体粒子物理等内容。
\subsection{历史}
天文学是一门古老的科学，长期以来与地球物理学相分离。在亚里士多德的宇宙观中，天体被认为是永恒不变的球体，其唯一的运动是匀速圆周运动；而地上世界则是生长与腐朽的领域，其中的自然运动沿直线进行，并在物体到达其目标时终止。因此，人们认为天界由一种与地球物质根本不同的物质构成——柏拉图认为那是火（Fire），而亚里士多德则称之为以太（Aether）。\cite{ref6,ref7}到了十七世纪，伽利略（Galileo）、笛卡尔（Descartes）和牛顿（Newton）等自然哲学家开始主张，天体与地上物体由相似的物质组成，并遵循相同的自然规律。\cite{ref8,ref9,ref10,ref11} 然而，他们所面临的困难在于：当时尚未发明出能够证明这些观点的科学工具。\cite{ref12}


在十九世纪的大部分时间里，天文学研究的重点仍是例行性的观测工作，如测量天体的位置与计算其运动。\cite{ref13,ref14}一种新的天文学——很快被称为天体物理学（Astrophysics）——开始兴起。威廉·海德·沃拉斯顿（William Hyde Wollaston）与约瑟夫·冯·夫琅和费（Joseph von Fraunhofer）分别独立发现，当太阳光被分解时，其光谱中出现了大量暗线（即光强减弱或缺失的区域）。\cite{ref15}

\paragraph{}
到 1860 年，物理学家古斯塔夫·基尔霍夫（Gustav Kirchhoff）与化学家罗伯特·本生（Robert Bunsen）证明，这些太阳光谱中的暗线与已知气体光谱中的亮线相对应，而每条特定的谱线都对应一种独特的化学元素。\cite{ref16}  
基尔霍夫推断，这些暗线是由太阳大气中化学元素吸收光线所造成的。\cite{ref17}  
由此，人们首次证明：太阳与恒星中存在的化学元素，也同样存在于地球上。

\paragraph{}
在进一步研究太阳与恒星光谱的人物中，诺曼·洛基尔（Norman Lockyer）尤为突出。1868 年，他在太阳光谱中发现了辐射线与暗线。与化学家爱德华·弗兰克兰（Edward Frankland）合作研究不同温度与压力下元素的光谱时，他注意到太阳光谱中一条黄色谱线无法对应任何已知元素。于是他提出，这条谱线代表一种新元素，并将其命名为“氦”（Helium），来源于希腊语“Helios”，意为“太阳之神”。\cite{ref18,ref19}

\paragraph{}
1885 年，爱德华·C·皮克林（Edward C. Pickering）在哈佛学院天文台发起了一项雄心勃勃的恒星光谱分类计划。他组织了一支女性计算员团队，其中包括威廉米娜·弗莱明（Williamina Fleming）、安东尼娅·莫里（Antonia Maury）与安妮·跳跃·坎农（Annie Jump Cannon）。她们根据照相底片上记录的光谱对恒星进行了系统分类。  
到 1890 年，团队已完成超过一万颗恒星的目录，并将它们划分为十三种光谱类型。根据皮克林的设想，到 1924 年，坎农将该目录扩展为九卷本、超过二十五万颗恒星，并发展出“哈佛分类系统”（Harvard Classification Scheme），该系统于 1922 年被国际天文学界正式采纳为全球标准。\cite{ref20}
